\documentclass[11pt]{article}

\usepackage[margin=1in]{geometry}
\usepackage{amsmath,amssymb,amsthm}
\usepackage{graphicx}

\newtheorem{theorem}{Theorem}
\newtheorem{lemma}{Lemma}
\renewcommand{\qedsymbol}{$\blacksquare$}

\title{MATH 327 Homework 1 Fall 2026}
\author{Yazan Ismail}
\date{10/7/2026}

\begin{document}

\maketitle

\section*{1. Problem 1.3}

\begin{lemma}\label{lem:sum}
For all \(n\in\mathbb{N}\),
\[
1+2+\cdots+n=\frac{n(n+1)}{2}.
\]
\end{lemma}

\begin{proof}
Let \(P(n)\) be the statement
\[
1+2+\cdots+n=\frac{n(n+1)}{2}.
\]
We proceed by induction on \(n\).

\textbf{Base case.} For \(n=1\),
\[
1=\frac{1(1+1)}{2},
\]
so \(P(1)\) holds.

\textbf{Inductive step.} Suppose \(P(m)\) holds for some \(m\in\mathbb{N}\).
Thus,
\begin{equation}\label{eq:sum-induction-hypothesis}
1+2+\cdots+m=\frac{m(m+1)}{2}.
\end{equation}
Adding \(m+1\) to both sides of \eqref{eq:sum-induction-hypothesis}, we get
\begin{align*}
1+2+\cdots+m+(m+1)
&=\frac{m(m+1)}{2}+(m+1)\\
&=(m+1)\left(\frac{m}{2}+1\right)\\
&=(m+1)\left(\frac{m+2}{2}\right)\\
&=\frac{(m+1)(m+2)}{2}.
\end{align*}
Therefore, \(P(m+1)\) holds, and the result follows by induction.
\end{proof}

\begin{theorem}
For all \(n\in\mathbb{N}\),
\[
1^3+2^3+\cdots+n^3=(1+2+\cdots+n)^2.
\]
\end{theorem}

\begin{proof}
Let \(P(n)\) be the statement
\[
1^3+2^3+\cdots+n^3=(1+2+\cdots+n)^2.
\]
We proceed by induction on \(n\).

\textbf{Base case.} For \(n=1\),
\[
1^3=1^2,
\]
so \(P(1)\) holds.

\textbf{Inductive step.} Suppose \(P(m)\) holds for some \(m\in\mathbb{N}\).
Thus,
\begin{equation}\label{eq:cubes-induction-hypothesis}
1^3+2^3+\cdots+m^3=(1+2+\cdots+m)^2.
\end{equation}
Adding \((m+1)^3\) to both sides of \eqref{eq:cubes-induction-hypothesis}, we get
\begin{align*}
1^3+2^3+\cdots+m^3+(m+1)^3
&=(1+2+\cdots+m)^2+(m+1)^3\\
&=\left(\frac{m(m+1)}{2}\right)^2+(m+1)^3
    &&\text{by Lemma~\ref{lem:sum}}\\
&=(m+1)^2\left(\frac{m^2}{4}+m+1\right)\\
&=(m+1)^2\left(\frac{m^2+4m+4}{4}\right)\\
&=(m+1)^2\left(\frac{m+2}{2}\right)^2\\
&=\left(\frac{(m+1)(m+2)}{2}\right)^2\\
&=(1+2+\cdots+m+(m+1))^2
    &&\text{by Lemma~\ref{lem:sum}}.
\end{align*}
Therefore, \(P(m+1)\) holds, and the result follows by induction.
\end{proof}

\section*{2. Problem 1.12(c)}

The following supporting identity is part 1.12(b).

\begin{lemma}[Pascal identity from 1.12(b)]
For \(n \in \mathbb{N}\) and \(k=1,2,\ldots,n\),
\[
\binom{n}{k}+\binom{n}{k-1}=\binom{n+1}{k}.
\]
\end{lemma}

\begin{proof}
Using the definition of the binomial coefficients,
\[
\binom{n}{k}=\frac{n!}{k!(n-k)!},
\qquad
\binom{n}{k-1}=\frac{n!}{(k-1)!(n-k+1)!},
\]
and
\[
\binom{n+1}{k}=\frac{(n+1)!}{k!(n-k+1)!}.
\]
Therefore,
\begin{align*}
\binom{n}{k}+\binom{n}{k-1}
&=\frac{n!}{k!(n-k)!}
  +\frac{n!}{(k-1)!(n-k+1)!}\\[4pt]
&=n!\left(
  \frac{n-k+1}{k!(n-k+1)!}
  +\frac{k}{k!(n-k+1)!}
  \right)\\[4pt]
&=n!\left(\frac{n+1}{k!(n-k+1)!}\right)\\[4pt]
&=\frac{(n+1)!}{k!(n-k+1)!}\\[4pt]
&=\binom{n+1}{k}.
\end{align*}
\end{proof}

\begin{theorem}[Binomial theorem]
For all \(a,b\in\mathbb{R}\) and \(n\in\mathbb{N}\),
\[
(a+b)^n=\sum_{k=0}^{n}\binom{n}{k}a^{n-k}b^k.
\]
\end{theorem}

\begin{proof}
Let \(P(n)\) be the stated binomial identity.
We proceed by induction on \(n\).

\textbf{Base case.} Let \(n=1\).
Then \(P(1)\) states
\begin{align*}
(a+b)^1
&=\binom{1}{0}a+\binom{1}{1}b\\
&=a+b.
\end{align*}

\textbf{Inductive step.} Suppose \(P(m)\) holds for some \(m\in\mathbb{N}\).
Thus, we have
\begin{align*}
(a+b)^m
&=\binom{m}{0}a^m+\binom{m}{1}a^{m-1}b\\
&\quad+\binom{m}{2}a^{m-2}b^2+\cdots\\
&\quad+\binom{m}{m-1}ab^{m-1}+\binom{m}{m}b^m\\
&=a^m+ma^{m-1}b+\frac12m(m-1)a^{m-2}b^2\\
&\quad+\cdots+mab^{m-1}+b^m.
\end{align*}
In these expanded expressions, only terms with valid indices are present when \(m\) is small.

\begin{align*}
(a+b)^{m+1}
&=(a+b)^m(a+b)\\[4pt]
&=(a+b)\Bigl[
   \binom{m}{0}a^m+\binom{m}{1}a^{m-1}b\\
&\qquad+\binom{m}{2}a^{m-2}b^2+\cdots\\
&\qquad+\binom{m}{m-1}ab^{m-1}+\binom{m}{m}b^m
   \Bigr]\\[6pt]
&=\binom{m}{0}a^{m+1}+\binom{m}{1}a^mb\\
&\quad+\binom{m}{2}a^{m-1}b^2+\cdots\\
&\quad+\binom{m}{m-1}a^2b^{m-1}+\binom{m}{m}ab^m\\
&\quad+\binom{m}{0}a^mb+\binom{m}{1}a^{m-1}b^2\\
&\quad+\binom{m}{2}a^{m-2}b^3+\cdots\\
&\quad+\binom{m}{m-1}ab^m+\binom{m}{m}b^{m+1}\\[6pt]
&=\binom{m}{0}a^{m+1}\\
&\quad+\left(\binom{m}{1}+\binom{m}{0}\right)a^mb\\
&\quad+\left(\binom{m}{2}+\binom{m}{1}\right)a^{m-1}b^2\\
&\quad+\cdots\\
&\quad+\left(\binom{m}{m}+\binom{m}{m-1}\right)ab^m\\
&\quad+\binom{m}{m}b^{m+1}.
\end{align*}
By part (b),
\[
\binom{m}{k}+\binom{m}{k-1}=\binom{m+1}{k}.
\]
Also,
\begin{align*}
\binom{m}{0}
&=\frac{m!}{0!(m-0)!}=1=\binom{m+1}{0},\\[6pt]
\binom{m}{m}
&=\frac{m!}{m!(m-m)!}=1=\binom{m+1}{m+1}.
\end{align*}
Hence,
\begin{align*}
(a+b)^{m+1}
&=\binom{m+1}{0}a^{m+1}\\
&\quad+\binom{m+1}{1}a^mb\\
&\quad+\binom{m+1}{2}a^{m-1}b^2+\cdots\\
&\quad+\binom{m+1}{m}ab^m\\
&\quad+\binom{m+1}{m+1}b^{m+1}\\[6pt]
&=a^{m+1}+(m+1)a^mb\\
&\quad+\frac12(m+1)m a^{m-1}b^2\\
&\quad+\cdots+(m+1)ab^m+b^{m+1}.
\end{align*}
Thus, \(P(m)\) implies \(P(m+1)\).
By induction, \(P(n)\) is true for all \(n\in\mathbb{N}\), so the binomial theorem holds.
\end{proof}

\section*{3. Bernoulli's inequality}

\begin{theorem}
For every \(x\geq -1\) and every \(n\in\mathbb{N}\),
\[
(1+x)^n\geq 1+nx.
\]
\end{theorem}

\begin{proof}
Fix \(x\in[-1,\infty)\), and let \(P(n)\) be the statement
\[
(1+x)^n\geq 1+nx.
\]
We prove \(P(n)\) for all \(n\in\mathbb{N}\) by induction.

\textbf{Base case.} For \(n=1\),
\[
(1+x)^1=1+x=1+1x,
\]
so (P(1)) holds.

\textbf{Inductive step.} Suppose \(P(m)\) holds for some \(m\in\mathbb{N}\).
Thus,
\[
(1+x)^m\geq 1+mx.
\]
Since \(x\geq -1\), we have \(1+x\geq 0\).
Multiplying both sides by (1+x) preserves the inequality:
\begin{align*}
(1+x)^{m+1}
&\geq(1+mx)(1+x)\\
&=1+(m+1)x+mx^2.
\end{align*}
Because \(m\in\mathbb{N}\) and \(x^2\geq 0\), we have \(mx^2\geq 0\).
Hence,
\[
(1+x)^{m+1}\geq 1+(m+1)x.
\]
Thus (P(m+1)) holds, and the result follows by induction.
\end{proof}

\section*{4. Problem 2.4}

\begin{lemma}\label{lem:sqrt-three}
\(\sqrt{3}\) is irrational.
\end{lemma}

\begin{proof}
Suppose, for the sake of contradiction, that \(\sqrt{3}\) is rational.
Then
\[
\sqrt{3}=\frac{p}{q}
\]
for some \(p\in\mathbb{Z}\) and \(q\in\mathbb{N}\).
We may assume that \(p\) and \(q\) are relatively prime.
Squaring both sides gives
\[
3q^2=p^2.
\]
We consider the possible parities of \(p\) and \(q\).

\textbf{Case 1: \(p\) and \(q\) are both odd.}
Write \(p=2k+1\) and \(q=2j+1\) for some \(k,j\in\mathbb{Z}\).
Then
\[
3(2j+1)^2=(2k+1)^2.
\]
Expanding and simplifying,
\begin{align*}
12j^2+12j+3&=4k^2+4k+1,\\
6j^2+6j+1&=2k^2+2k.
\end{align*}
The left side of the last equality is odd, while the right side is even, a contradiction.

\textbf{Case 2: \(p\) and \(q\) have opposite parity.}
Then \(p^2\) and \(q^2\) have opposite parity.
Since \(3q^2\) has the same parity as \(q^2\), it cannot equal \(p^2\), a contradiction.

\textbf{Case 3: \(p\) and \(q\) are both even.}
Then \(p\) and \(q\) have a common factor of \(2\), contradicting the assumption that they are relatively prime.

These cases exhaust all possible parities of \(p\) and \(q\).
Therefore, \(\sqrt{3}\) is irrational.
\end{proof}

\begin{theorem}
\(\sqrt[3]{5-\sqrt{3}}\) is irrational.
\end{theorem}

\begin{proof}
Suppose, for the sake of contradiction, that \(\sqrt[3]{5-\sqrt{3}}\) is rational.
Then
\[
\sqrt[3]{5-\sqrt{3}}=\frac{p}{q}
\]
for some \(p\in\mathbb{Z}\) and \(q\in\mathbb{N}\).
Cubing both sides gives
\[
5-\sqrt{3}=\frac{p^3}{q^3}.
\]
Rearranging,
\[
\sqrt{3}=\frac{5q^3-p^3}{q^3}.
\]
Since \(5q^3-p^3\in\mathbb{Z}\) and \(q^3\in\mathbb{N}\), this would make \(\sqrt{3}\) rational, contradicting Lemma~\ref{lem:sqrt-three}.
Therefore, \(\sqrt[3]{5-\sqrt{3}}\) is irrational.
\end{proof}

\section*{5. Problem 3.6(b)}

\begin{theorem}
For every \(n\in\mathbb{N}\) and all \(a_1,a_2,\ldots,a_n\in\mathbb{R}\),
\[
\left|a_1+a_2+\cdots+a_n\right|
\leq |a_1|+|a_2|+\cdots+|a_n|.
\]
\end{theorem}

\begin{proof}
Let \(P(n)\) be the statement that, for every \(a_1,a_2,\ldots,a_n\in\mathbb{R}\),
\[
\left|a_1+a_2+\cdots+a_n\right|
\leq |a_1|+|a_2|+\cdots+|a_n|.
\]
We proceed by induction on \(n\).

\textbf{Base case.} For \(n=1\) and any \(a_1\in\mathbb{R}\),
\[
|a_1|\leq |a_1|.
\]
Thus \(P(1)\) holds.

\textbf{Inductive step.} Suppose \(P(m)\) holds for some \(m\in\mathbb{N}\).
Let \(a_1,a_2,\ldots,a_{m+1}\in\mathbb{R}\).
By the induction hypothesis,
\begin{equation}\label{eq:triangle-induction-hypothesis}
\left|a_1+a_2+\cdots+a_m\right|
\leq |a_1|+|a_2|+\cdots+|a_m|.
\end{equation}
Adding \(|a_{m+1}|\) to both sides of \eqref{eq:triangle-induction-hypothesis} gives
\[
\left|a_1+a_2+\cdots+a_m\right|+|a_{m+1}|
\leq |a_1|+|a_2|+\cdots+|a_m|+|a_{m+1}|.
\]
By the triangle inequality,
\[
\left|(a_1+a_2+\cdots+a_m)+a_{m+1}\right|
\leq \left|a_1+a_2+\cdots+a_m\right|+|a_{m+1}|.
\]
Therefore, by transitivity,
\[
\left|a_1+a_2+\cdots+a_m+a_{m+1}\right|
\leq |a_1|+|a_2|+\cdots+|a_m|+|a_{m+1}|.
\]
Thus \(P(m+1)\) holds, and the result follows by induction.
\end{proof}

\section*{6. Bounds table for Question 4.1}

The upper- and lower-bound entries below give all real bounds for each set.
\begin{center}
\renewcommand{\arraystretch}{1.7}
\resizebox{\textwidth}{!}{%
\begin{tabular}{|c|c|c|c|c|c|c|}
\hline
Set & Upper bounds & Lower bounds & Maximum & Minimum & Supremum & Infimum \\
\hline
(h) & \(\text{DNE}\) & \(\{x\in\mathbb{R}:x\leq 2\}\)
    & \(\text{DNE}\) & \(2\) & \(\text{DNE}\) & \(2\) \\
\hline
(i) & \(\{x\in\mathbb{R}:x\geq 1\}\) & \(\{x\in\mathbb{R}:x\leq 0\}\)
    & \(1\) & \(0\) & \(1\) & \(0\) \\
\hline
(j) & \(\{x\in\mathbb{R}:x\geq 1\}\) & \(\{x\in\mathbb{R}:x\leq \frac{2}{3}\}\)
    & \(\text{DNE}\) & \(\frac{2}{3}\) & \(1\) & \(\frac{2}{3}\) \\
\hline
(k) & \(\text{DNE}\) & \(\{x\in\mathbb{R}:x\leq 0\}\)
    & \(\text{DNE}\) & \(0\) & \(\text{DNE}\) & \(0\) \\
\hline
(n) & \(\{x\in\mathbb{R}:x\geq\sqrt{2}\}\) & \(\{x\in\mathbb{R}:x\leq-\sqrt{2}\}\)
    & \(\text{DNE}\) & \(\text{DNE}\) & \(\sqrt{2}\) & \(-\sqrt{2}\) \\
\hline
(r) & \(\{x\in\mathbb{R}:x\geq 1\}\) & \(\{x\in\mathbb{R}:x\leq 1\}\)
    & \(1\) & \(1\) & \(1\) & \(1\) \\
\hline
(w) & \(\{x\in\mathbb{R}:x\geq\frac{\sqrt{3}}{2}\}\)
    & \(\{x\in\mathbb{R}:x\leq-\frac{\sqrt{3}}{2}\}\)
    & \(\frac{\sqrt{3}}{2}\) & \(-\frac{\sqrt{3}}{2}\)
    & \(\frac{\sqrt{3}}{2}\) & \(-\frac{\sqrt{3}}{2}\) \\
\hline
\end{tabular}%
}
\end{center}

\section*{7. Problem 4.7}

\begin{theorem}[Part (a)]\label{thm:subset-bounds}
Let \(S\) and \(T\) be nonempty bounded subsets of \(\mathbb{R}\).
If \(S\subseteq T\), then
\[
\inf T\leq\inf S\leq\sup S\leq\sup T.
\]
\end{theorem}

\begin{proof}
Suppose \(S\subseteq T\).

\noindent\underline{\textbf{Supremum case.}}
Suppose, for the sake of contradiction, that \(\sup S>\sup T\).
Since \(\sup T\) is an upper bound for \(T\), every \(s\in S\) also satisfies \(s\leq\sup T\).
Thus \(\sup T\) is an upper bound for \(S\).
But \(\sup T<\sup S\), contradicting the fact that \(\sup S\) is the least upper bound for \(S\).
Hence \(\sup S\leq\sup T\).

\noindent\underline{\textbf{Infimum case.}}
Suppose, for the sake of contradiction, that \(\inf S<\inf T\).
Since \(\inf T\) is a lower bound for \(T\), every \(s\in S\) also satisfies \(\inf T\leq s\).
Thus \(\inf T\) is a lower bound for \(S\).
But \(\inf T>\inf S\), contradicting the fact that \(\inf S\) is the greatest lower bound for \(S\).
Hence \(\inf T\leq\inf S\).

Since \(S\) is nonempty, choose \(s\in S\).
The definitions of infimum and supremum give \(\inf S\leq s\leq\sup S\), so \(\inf S\leq\sup S\).
Combining the three inequalities yields
\[
\inf T\leq\inf S\leq\sup S\leq\sup T.
\]
\end{proof}

\begin{theorem}[Part (b)]
Let \(S\) and \(T\) be nonempty bounded subsets of \(\mathbb{R}\).
Then
\[
\sup(S\cup T)=\max\{\sup S,\sup T\}.
\]
\end{theorem}

\begin{proof}
Let \(x=\max\{\sup S,\sup T\}\) and \(\alpha=\sup(S\cup T)\).
Since \(S\subseteq S\cup T\) and \(T\subseteq S\cup T\), part (a) gives
\[
\sup S\leq\alpha
\qquad\text{and}\qquad
\sup T\leq\alpha.
\]
Therefore \(x\leq\alpha\).

Suppose, for the sake of contradiction, that \(x<\alpha\).
Then \(0<\alpha-x\), so
\[
0<\frac{\alpha-x}{2}<\alpha-x.
\]
Adding \(x\) gives
\[
x<x+\frac{\alpha-x}{2}<\alpha.
\]
Since \(x\geq\sup S\) and \(x\geq\sup T\), \(x\) is an upper bound for both \(S\) and \(T\).
Consequently, \(x+\frac{\alpha-x}{2}\) is an upper bound for \(S\cup T\).
This contradicts the fact that \(\alpha\) is the least upper bound for \(S\cup T\).
Thus \(x=\alpha\), or
\[
\sup(S\cup T)=\max\{\sup S,\sup T\}.
\]
\end{proof}

\section*{8. Problem 4.8(b)}

\begin{theorem}
Let \(S\) and \(T\) be nonempty subsets of \(\mathbb{R}\) such that
\[
s\leq t
\qquad\text{for all }s\in S\text{ and }t\in T.
\]
Then \(\sup S\leq\inf T\).
\end{theorem}

\begin{proof}
Let \(S\) and \(T\) be nonempty subsets of \(\mathbb{R}\) with the following property:
\[
s\leq t \quad\text{for all }s\in S\text{ and }t\in T.
\]

Suppose for the sake of contradiction that \(\sup(S)>\inf(T)\).
Thus, \(\inf(T)\) is not an upper bound for \(S\).
Hence, there exists an \(s\in S\) such that \(\inf(T)<s<\sup(S)\).
Let \(s_1\) be such \(s\).
Now we have two cases:

\textbf{Case 1.} \(s_1\) is a lower bound for \(T\), in which case, \(s_1\) is a greater lower bound for \(T\) than \(\inf(T)\), which is a contradiction.

\textbf{Case 2.} \(s_1\in T\).
Since \(s_1\leq t\) for all \(s\in S\) and \(t\in T\), then \(\min(T)=s_1\).
Since \(s_1\) is the minimum of \(T\), then \(s_1=\inf(T)\), contradicting \(\inf(T)<s_1\).

\textbf{Case 3.} \(s_1\) is an upper bound for \(T\), in which case, for all \(t\in T\), \(t\leq s_1\).
If \(t<s_1\) for any \(t\in T\), then we contradict \(s\leq t\) for all \(s\in S\) and \(t\in T\).
Thus, \(t=s_1\) for all \(t\in T\).
However, since \(s_1<\sup(S)\), then there exists \(s_2\in S\) such that \(s_1<s_2<\sup(S)\).
Thus, \(s_2\in S\) such that \(s_2>t\) for all \(t\in T\), contradicting our supposition that for all \(s\in S\) and for all \(t\in T\), \(s\leq t\).

Hence, we've exhausted all cases with contradictions.

Therefore, \(\sup(S)\leq\inf(T)\).
\end{proof}

\end{document}
